% =============================================================================
% Project : THE_PHILOSOPHY_OF_VIABILITY
% Book    : Book I — Ontology of Viability
% Book ID : VS.B01
%
% File    :
%   2_MAINMATTER/PART_I_IDENTITY_AND_MANIFESTATION/
%   CHAPTER_02_INVARIANCE/sections/
%   section_04_invariance_and_transformation.tex
%
% Section : 04 — Invariance and Transformation
% Chapter : 02 — Invariance
%
% Purpose :
%   To examine the ontological and analytical relation between
%   transformation, preservation, and alteration, establishing
%   the conditions under which invariance may be attributed to
%   a specified aspect of an existent undergoing transformation.
%
% Conceptual Scope :
%   Transformation, preservation, alteration, partial preservation,
%   total preservation, configurations of determinations,
%   relational organization, and the limits of invariance claims.
%
% =============================================================================

\section{Invariance and Transformation}
\label{sec:invariance-and-transformation}

\subsection{The Relation Between Invariance and Transformation}

Invariance and transformation are sometimes treated as
opposing conditions. According to this interpretation,
whatever remains invariant must remain unchanged, whereas
whatever undergoes transformation must necessarily lose
its previous determinations.

This opposition is inadequate because it treats change
and preservation as properties that must characterize
an existent uniformly and in its entirety.

A transformation need not affect every determination,
relation, or organizational feature of an existent in
the same manner. Some aspects may be altered, while others
may remain preserved relative to the transformation
under consideration.

Conversely, the preservation of a particular determination
or relation does not establish that the existent has
remained unchanged in every other respect.

Transformation and invariance therefore belong to
analytically distinguishable but potentially compatible
descriptions of an existent and its determinations.

Their relation cannot be adequately examined through
the undifferentiated opposition between change and
unchange. It requires identifying what is transformed,
what is assessed as preserved, and the respect in which
preservation or alteration is attributed.

The central question is consequently not whether an
existent is either transformed or invariant without
qualification. It is whether a specified feature,
determination, or relation remains preserved under
a specified transformation, within a specified domain,
and according to an appropriate criterion.

Invariance does not negate transformation, and
transformation does not, merely by occurring, negate
every possibility of invariance.

\subsection{Transformation Does Not Necessarily Eliminate Invariance}

A transformation involves an alteration in at least
one relevant aspect of a configuration, determination,
or relation, such that the resulting condition differs
from the prior condition in the respect under examination.

However, the occurrence of such an alteration does not
establish that every other aspect has also been altered.

A transformation may affect certain determinations
while leaving others preserved. It may modify quantitative
properties without eliminating a specified structural
relation. It may alter the properties of relata while
leaving a particular relational organization intact.

In these cases, the transformation and the preservation
of a specified feature are compatible.

The preservation claim must nevertheless identify
the feature whose invariance is being assessed. It is
not sufficient to observe that something appears to
have remained unchanged. The relevant determination
or relation must be specified, and its preservation
must be evaluated under the transformation in question.

The distinction is important because the existence
of alteration in one respect provides no sufficient
grounds for denying preservation in another.

A transformation that alters the quantitative properties
of certain constituents, for example, does not thereby
establish that every relation among those constituents
has also been altered.

Likewise, a transformation that changes the material
constitution of an existent does not, without further
analysis, determine whether its structural organization
has been preserved or altered.

The possibility of invariance within transformation
therefore depends on the differentiation of the aspects
under examination. It does not require that the existent
remain wholly unaffected.

\subsection{Invariance Does Not Imply the Absence of Transformation}

The converse limitation is equally necessary.

The preservation of a determination or relation does
not establish that no transformation has occurred.

An invariant may be identified within a configuration
that has undergone alteration in other respects.
The preservation of one feature cannot be generalized
to the entire existent unless the scope of the claim
explicitly warrants such a generalization.

For example, the preservation of a specified relational
pattern does not establish that every participating
determination has remained unchanged. Similarly, the
preservation of a structural feature does not establish
that the structure has remained identical in every
respect or that no other aspect of its organization
has been modified.

The error arises when a limited preservation claim
is extended beyond its subject and criterion.

An invariance claim concerns the preservation of
a specified feature under specified conditions. It
does not automatically characterize the total condition
of the existent.

Accordingly, the assertion that a particular determination
remains invariant must not be converted into the assertion
that the existent as a whole has undergone no transformation.

The presence of an invariant therefore does not establish
immutability. It establishes, at most, the preservation
of the specified feature within the scope of the
transformation and criterion under consideration.

\subsection{The Coexistence of Preservation and Alteration}

Preservation and alteration may coexist within the
same transformation.

This possibility arises because the determinations
and relations associated with an existent need not
be affected uniformly.

A transformation may preserve one determination,
alter another, modify a relation, and establish a
different configuration among other determinations.

These outcomes are not mutually exclusive merely
because they occur within the same transformed
configuration.

Their compatibility depends on the respects in
which preservation and alteration are attributed.

A contradiction would arise if the same subject,
the same determination or relation, and the same
respect were asserted to be both preserved and altered
under the same transformation, within the same domain,
and according to the same criterion.

Where the relevant subjects or respects differ,
the two attributions may consistently coexist.

For example, a particular structural relation may
remain preserved while the quantitative properties
of the elements participating in that relation are
altered. The relational feature and the quantitative
properties are not the same subject of assessment.
Their different outcomes therefore do not constitute
a contradiction.

Likewise, the preservation of an organizational
pattern may coexist with the alteration of the
determinations occupying particular positions within
that pattern.

The distinction does not mean that every apparent
coexistence of preservation and alteration is coherent.
The subject, transformation, domain, respect, and
criterion must be sufficiently specified to determine
whether the two claims concern compatible aspects
or make incompatible assertions about the same one.

The coexistence of preservation and alteration is
therefore a possibility that requires analytical
specification, not a conclusion that can be assumed
in every case.

\subsection{Partial and Total Preservation}

The distinction between partial and total preservation
concerns the extent of what is preserved relative to
a specified assessment.

Partial preservation occurs when some of the features
included in the assessment remain preserved while
others are altered.

It does not necessarily indicate that the transformation
is minor, nor does it determine the ontological
significance of the features that remain preserved.
A transformation may be extensive in several respects
while preserving a particular determination or relation.

Total preservation requires a more restricted and
carefully specified meaning.

It can be attributed only when every feature included
within the explicitly defined scope of assessment
remains preserved according to the relevant criterion.

The qualification concerning scope is indispensable.
No finite or restricted assessment can establish
that every aspect of an existent, without qualification,
has remained unchanged.

Consequently, total preservation within a specified
assessment must not be confused with absolute
immutability or the absence of transformation in
every conceivable respect.

Even when all features included in an assessment
remain preserved, the assessment may exclude other
determinations, relations, or aspects of the existent.
The conclusion must therefore remain proportionate
to the scope of the examination.

Partial and total preservation are thus comparative
descriptions of the extent to which specified features
are preserved. Neither description, by itself,
establishes the numerical identity of the existent.

Nor does the distinction determine whether the
preserved features are ontologically fundamental,
structurally significant, or merely restricted
characteristics. Such judgments require additional
arguments and cannot be derived from the extent
of preservation alone.

\subsection{Transformation and the Configuration of Determinations}

Transformation may affect individual determinations
as well as the configuration in which those
determinations are organized.

These two levels of assessment must remain distinct.

A determination may remain preserved while its
position or role within a configuration changes.
Conversely, a configuration may retain a specified
organizational relation even when some of the
determinations participating in it undergo alteration.

The preservation of an individual determination
therefore does not automatically establish the
preservation of the configuration in which it occurs.

Similarly, the preservation of an organizational
relation does not require that every determination
participating in that relation remain unchanged.

The distinction becomes especially important when
the configuration itself is the subject of an
invariance claim.

In such cases, the analysis must identify whether
the relevant feature concerns:

\begin{enumerate}
    \item the preservation of particular determinations;
    \item the preservation of relations among determinations;
    \item the preservation of a specified organization
    of determinations and relations; or
    \item the preservation of a combination of these
    features under the transformation in question.
\end{enumerate}

These are distinct possible subjects of assessment.
They must not be treated as interchangeable.

A configuration may undergo alteration even when
some of its constituent determinations remain
preserved. It may also retain a specified organizational
feature despite changes in some of its constituents.

Furthermore, the preservation of a relation does
not necessarily entail the preservation of the
determinations participating in it. Whether a
particular relation can remain preserved under
such changes depends on the nature of the relation
and the criterion by which its preservation is assessed.

The analysis of invariance must therefore distinguish
between the preservation of determinations, the
preservation of relations, and the preservation
of configurations.

No one of these forms of preservation should be
assumed to establish the others without further
justification.

\subsection{Transformation and the Limits of Generalization}

An invariant established under a particular
transformation cannot automatically be generalized
to other transformations.

A determination or relation may remain preserved
under one transformation while undergoing alteration
under another. The same feature may also be preserved
within a restricted domain but fail to remain
preserved when the transformation extends beyond
that domain.

The validity of an invariance claim must therefore
remain bounded by the conditions under which it
has been established.

These conditions include the specified subject,
the transformation or admissible class of transformations,
the domain of applicability, the respect of preservation,
and the criterion used to assess the result.

This limitation does not make invariance arbitrary.
It establishes the conditions under which a claim
is warranted and prevents a restricted instance
of preservation from being treated as an unrestricted
ontological principle.

A feature demonstrated to remain preserved under
a defined class of transformations may legitimately
be described as invariant relative to that class,
provided that the relevant conditions and limits
have been established.

Such a conclusion does not establish that the
feature is preserved under every possible
transformation, nor does it establish that the
same result holds in every domain or respect.

The distinction between restricted and unrestricted
claims is therefore essential.

An unrestricted claim would require a correspondingly
unrestricted justification. It cannot be inferred
merely from the successful establishment of invariance
in a limited set of cases.

The scope of an invariance claim must consequently
remain proportionate to the scope of the argument
supporting it.

\subsection{Transformation, Invariance, and the Question of Identity}

The coexistence of preservation and alteration
raises a further question concerning the relation
between transformation and identity.

The preservation of a determination or relation
across a transformation may contribute to the
characterization of the existent and the manner
in which its configuration has changed.

Nevertheless, the existence of an invariant does
not, by itself, establish that the existent considered
before and after the transformation is numerically
identical.

An invariant specifies a feature that remains
preserved under defined conditions. Numerical
identity concerns a different ontological question.
The former cannot be substituted for an argument
establishing the latter.

The converse limitation is equally important.

The alteration or disappearance of a previously
identified invariant does not, by itself, establish
the disappearance of the identity of the existent.

The loss of a particular determination, relation,
or organizational feature may be relevant to an
investigation of identity, but its ontological
significance cannot be determined solely by the
fact that the feature has ceased to be preserved.

These limitations follow from the analytical
distinction between identity and invariance.
Neither concept can be reduced to the other, and
neither can independently settle every question
concerning the other.

Transformation may therefore be examined in terms
of preserved and altered features without presuming
that these findings resolve the identity of the
existent.

The question of what remains invariant and the
question of whether the existent remains numerically
identical must remain distinct, even when their
investigation is related.

\subsection{Analytical Possibilities of Preservation and Alteration}

The relation between transformation, preservation,
and alteration can be clarified through several
analytical possibilities.

A transformation may involve the preservation of
a specified determination or relation while other
aspects remain outside the scope of assessment.

It may involve the alteration of a specified
determination or relation without any invariant
having been established.

It may also involve the coexistence of preservation
and alteration, provided that the respective claims
concern compatible subjects or respects.

These possibilities clarify the logical relations
between transformation and invariance. They do
not constitute an exhaustive classification of
transformations, nor do they imply that every
transformation must produce an identifiable invariant.

The analytical distinctions can be stated as follows:

\begin{enumerate}
    \item \textbf{Preservation under transformation:}
    A specified feature remains preserved according
    to a defined criterion, notwithstanding alteration
    elsewhere in the configuration.

    \item \textbf{Alteration under transformation:}
    A specified feature undergoes alteration. This
    finding does not, by itself, determine whether
    other features remain preserved.

    \item \textbf{Coexisting preservation and alteration:}
    Some specified features remain preserved while
    others undergo alteration within the same
    transformation.

    \item \textbf{Undetermined preservation:}
    A transformation is established, but the available
    analysis does not determine whether a particular
    feature has been preserved, altered, or replaced.
\end{enumerate}

The fourth possibility is important because the
absence of an established invariant must not be
confused with the establishment of alteration.
An insufficiently determined case remains
unresolved with respect to the relevant preservation
claim.

Likewise, the identification of alteration in one
respect does not establish alteration in every
other respect.

These possibilities therefore distinguish what
may be asserted from what remains undetermined.
They do not establish a universal sequence or
a necessary outcome for every transformation.

The analytical value of these distinctions lies
in preventing the reduction of transformation
to either complete preservation or complete
alteration.

\subsection{Synthesis}

Invariance and transformation are not mutually
exclusive conditions.

A transformation may preserve certain determinations,
relations, or organizational features while altering
others. The occurrence of alteration in one respect
does not necessarily eliminate invariance in another,
and the preservation of a particular feature does
not establish the absence of transformation in
the existent as a whole.

The distinction between partial and total preservation
clarifies the extent of a preservation claim.
Total preservation can be asserted only relative
to an explicitly specified scope of assessment;
it must not be confused with absolute immutability.

The analysis must also distinguish between the
preservation of individual determinations, relations
among determinations, and the organization of a
configuration. These forms of preservation are
related but not interchangeable.

Every invariance claim remains bounded by its
specified subject, transformation, domain, respect,
and criterion. Its conclusions cannot be generalized
beyond the conditions under which preservation
has been established.

Finally, neither the preservation nor the alteration
of particular features is sufficient, by itself,
to establish or negate the numerical identity
of an existent.

Transformation and invariance therefore provide
distinct but related perspectives for examining
what is preserved and what is altered. Their
analysis contributes to the characterization of
an existent without replacing the separate
ontological question of its identity.

The next section examines the relation between
invariance and difference, addressing how the
preservation of specified features may coexist
with difference within and across transformed
configurations.

% =============================================================================
% END OF SECTION 04 OF CHAPTER 02
% =============================================================================